% ============================================================
%  Meridian Capital - Asset Management Challenge 2026
%  Round 2 presentation (Beamer, 16:9).
%  Compile with XeLaTeX, twice. DM Sans font files are included.
%  Files needed: Meridian-Capital-Official-Logo.png (next to this file or in figures/)
% ============================================================
\documentclass[aspectratio=169,11pt,t]{beamer}
\usepackage{colortbl}

\usepackage{fontspec}
\setmainfont{DMSans-Regular.ttf}[BoldFont=DMSans-Bold.ttf,ItalicFont=DMSans-Italic.ttf]
\setsansfont{DMSans-Regular.ttf}[BoldFont=DMSans-Bold.ttf,ItalicFont=DMSans-Italic.ttf]
\usefonttheme{professionalfonts}
\usepackage{microtype}
\usepackage{booktabs,array,tabularx}
\usepackage{graphicx}
\graphicspath{{./}{figures/}}
\usepackage{tikz}

% ---------- palette ----------
\definecolor{navy}{HTML}{0A2233}
\definecolor{ink}{HTML}{173242}
\definecolor{gold}{HTML}{A59577}
\definecolor{paper}{HTML}{F7F5EF}
\definecolor{muted}{HTML}{59666B}
\definecolor{line}{HTML}{D9DCD7}

\newcommand{\hd}[1]{\textcolor{navy}{\textbf{#1}}}
\newcommand{\smnote}[1]{{\footnotesize\color{muted}#1}}

% ---------- beamer theme ----------
\usetheme{default}
\usecolortheme{default}
\setbeamertemplate{navigation symbols}{}
\setbeamercolor{normal text}{fg=ink,bg=white}
\setbeamercolor{structure}{fg=navy}
\setbeamercolor{frametitle}{fg=navy}
\setbeamercolor{itemize item}{fg=gold}
\setbeamercolor{itemize subitem}{fg=gold}
\setbeamercolor{enumerate item}{fg=gold}
\setbeamertemplate{itemize item}[square]
\setbeamertemplate{itemize subitem}[square]
\setbeamerfont{itemize/enumerate body}{size=\normalsize}
\setbeamerfont{itemize/enumerate subbody}{size=\small}

% body, title and footer share the same side margins
\setbeamersize{text margin left=1.2cm,text margin right=1.2cm}

% frame title: gold rule below
\setbeamertemplate{frametitle}{%
  \vspace{0.55em}
  \begin{beamercolorbox}[wd=\paperwidth,ht=2.0ex,dp=0.3ex,leftskip=1.2cm,rightskip=1.2cm]{frametitle}
    {\rmfamily\LARGE\color{navy}\insertframetitle}
  \end{beamercolorbox}
  \vspace{-1.0em}
  {\hspace{1.2cm}\color{gold}\rule{\dimexpr\paperwidth-2.4cm\relax}{0.8pt}}
}

% footer: MERIDIAN CAPITAL · AMC 2026, page number right
\setbeamertemplate{footline}{%
  \leavevmode\hbox{\begin{beamercolorbox}[wd=\paperwidth,ht=3ex,dp=1ex,leftskip=1.2cm,rightskip=1.2cm]{normal text}
    {\color{line}\hrule height 0.3pt}\vspace{2pt}
    {\fontsize{5}{6}\selectfont\color{muted}{\addfontfeatures{LetterSpace=7}MERIDIAN CAPITAL}\ ·\ AMC 2026}%
    \hfill{\fontsize{5}{6}\selectfont\color{navy}\insertframenumber\,/\,\inserttotalframenumber}
  \end{beamercolorbox}}%
}

% highlighted box
\usepackage[most]{tcolorbox}
\newtcolorbox{calloutbox}[1]{enhanced, colback=paper, colframe=paper,
  borderline west={2.5pt}{0pt}{gold}, boxrule=0pt, arc=0pt,
  left=10pt, right=8pt, top=6pt, bottom=6pt,
  title={#1}, fonttitle=\bfseries\small\raggedright, coltitle=navy, colbacktitle=paper,
  attach title to upper=\par\vspace{2pt}, before upper={\color{ink}\raggedright}}
\newenvironment{callout}[1]{\begin{calloutbox}{#1}}{\end{calloutbox}}

% big number tile: value, label, width
\newcommand{\kpi}[3]{%
  \begin{minipage}[t]{#3}\centering
  {\rmfamily\fontsize{22}{26}\selectfont\color{navy}#1}\\[3pt]
  {\scriptsize\color{gold}{\addfontfeatures{LetterSpace=5}\MakeUppercase{#2}}}
  \end{minipage}}

% compact list inside columns
\newcommand{\tight}{\setlength{\itemsep}{0.45em}\setlength{\parskip}{0pt}}

% ============================================================
\begin{document}

% ---------------------------- TITLE ----------------------------
{
\setbeamertemplate{footline}{}
\setbeamercolor{background canvas}{bg=paper}
\begin{frame}[plain]
\centering
\vspace*{0.1cm}
\includegraphics[width=5.2cm]{Meridian-Capital-Official-Logo.png}\par
\vspace{0.25cm}
{\color{gold}\rule{3cm}{0.8pt}}\par
\vspace{0.35cm}
{\rmfamily\fontsize{22}{26}\selectfont\color{navy}Tactical Deviation\par}
\vspace{0.15cm}
{\small\color{muted}Discretionary Equity Investing\par}
\vspace{0.15cm}
{\scriptsize\color{gold}{\addfontfeatures{LetterSpace=7}ROUND 2 \,·\, ASSET MANAGEMENT CHALLENGE 2026}\par}
\vfill
{\footnotesize\color{navy}Alessia Battistella \ · \ Tze Qing Ginnie Lee \ · \ Katherine Lerner-Teves\\[1pt] Kyla Kaplan \ · \ Lorenzo Comi\par}
\vspace{0.15cm}
{\scriptsize\color{muted}USI \,·\, November 2026\par}
\vspace{0.2cm}
\end{frame}
}

% ---------------------------- AGENDA ----------------------------
\begin{frame}{Agenda}
\vspace{1.2em}
\begin{enumerate}\setlength{\itemsep}{0.8em}
\item Performance of the first round
\item Tactical deviation and risk
\item Stock analysis and valuation
\item Alphas, portfolio weights and credibility
\end{enumerate}
\end{frame}

\begin{frame}{Performance of the First Round}
\vspace{.1em}
\begin{center}\small\begin{tabular}{lr}\toprule
16 October--6 November 2026 & Result\\\midrule
Return & $-2.885\%$\\
Ending NAV & \$97,115,068.58\\
Annualized realized volatility & 8.14\%\\
SAA volatility target & 8.00\%\\
Annualized realized Sharpe & $-6.50$\\\bottomrule
\end{tabular}\end{center}
\begin{itemize}\tight
\item NAV includes both VGSH and STIP distributions (\$100,916.79 in total) and opening cash.
\item Just 15 daily returns: annualized Sharpe is unstable and describes this short losing period.
\end{itemize}
\smnote{Public-close reconstruction of the opening book, not a Bloomberg-certified series. Volatility and Sharpe use daily returns and the archived Treasury proxy.}
\end{frame}
\begin{frame}{Performance Attribution}
\vspace{.1em}
\begin{center}\small\begin{tabular}{lrr}\toprule
Position & P\&L, USD & Return contribution, pp\\\midrule
SPY&-955,365&-0.955\\
EEMV&-346,050&-0.346\\
DBC&-1,757,263&-1.757\\
SCHD&78,300&+0.078\\
VGSH&31,106&+0.031\\
STIP&64,341&+0.064\\
\midrule Total & $-2,884,931$ & $-2.885$\\\bottomrule\end{tabular}\end{center}
\textbf{DBC drove about 61\% of the net loss.} SPY was the second-largest detractor; the other three positive holdings partly offset the losses.
\par\medskip\smnote{Round 2 reduces the commodity target from 26.95\% to 17.97\%, but DBC plus Exxon still leaves material energy exposure. Contributions include dividends and use initial \$100m NAV.}
\end{frame}
% ---------------------------- 1. TACTICAL DEVIATION ----------------------------
\begin{frame}{Funding the Deviation}
\vspace{.1em}
\begin{columns}[T]
\begin{column}{.50\textwidth}
\textbf{Same benchmark, new stock sleeve}
\begin{itemize}\tight
\item Original six-ETF SAA stays the benchmark.
\item One-third of post-loss NAV goes to six stocks, following the class instruction.
\item We scale each ETF target by two-thirds. This funding choice preserves relative ETF proportions.
\end{itemize}
\end{column}
\begin{column}{.44\textwidth}\small
\begin{tabular}{lrr}\toprule
ETF & SAA & New target\\\midrule
SPY&26.08\%&17.39\%\\
EEMV&9.66\%&6.44\%\\
DBC&26.95\%&17.97\%\\
SCHD&3.28\%&2.19\%\\
VGSH&5.00\%&3.33\%\\
STIP&29.03\%&19.35\%\\\bottomrule
\end{tabular}
\end{column}\end{columns}
\vspace{.8em}\smnote{NAV: \$97,115,068.58. Stock budget: \$32,371,689.53. No fees or reserve. Entered 7 November at 6 November closing prices. Whole shares leave \$116.05 cash.}
\end{frame}


\begin{frame}{Tracking Error and Risk}
\vspace{.1em}
\begin{center}\small
\begin{tabular}{lrr}\toprule
& SAA & New portfolio\\\midrule
Annual monthly tracking error&0.00\%&3.92\%\\
Annual monthly volatility&7.89\%&9.47\%\\
One-month 5\% VaR, loss / NAV&3.74\%&4.50\%\\\bottomrule
\end{tabular}\end{center}
\begin{itemize}\tight
\item Official TE uses 60 complete monthly USD returns: November 2021--October 2026. It passes the 2--5\% band.
\item Recent daily TE is higher: 6.53\% over 252 sessions and 7.74\% over 63. Our extra daily 5\% preference fails.
\item 217 post-DBC-change 33-session replays average $-0.395\%$ active return. Overlap leaves only 6--7 non-overlapping windows per offset.
\end{itemize}
\smnote{Zero-mean normal VaR = 1.645 $\times$ annual volatility / $\sqrt{12}$. Replays reuse today's selected universe and do not establish a future return advantage.}
\end{frame}


% ---------------------------- 2. STOCK ANALYSIS ----------------------------
% one frame per stock: copy the block 3 to 6 times
\begin{frame}{ADBE -- Adobe Inc.}
\vspace{.1em}
\textbf{Value:} \$18.145 FY2026 GAAP EPS $\times$ 16 = \$290.32, versus \$239.30 entry. FCFF cross-check: \$271.14.
\begin{itemize}\tight
\item Subscription cash generation supports the case. SBC is already in GAAP EPS and is charged in the cash model.
\item \textbf{Catalyst: 9 December}, confirmed Q4 results and FY2027 outlook. Watch renewals, paid AI revenue and pricing.
\item \textbf{Main risk:} cheaper AI tools or copied features weaken retention. A 20\% cash reduction values the stock at \$216.54.
\item \textbf{What would change our view:} guidance cuts or cash deterioration that confirms structural disruption.
\end{itemize}
\smnote{Fund target 4.65\%. Original annual alpha 12.07\%. Sources: Adobe Q3 FY2026 release and investor calendar.}
\end{frame}


\begin{frame}{KR -- The Kroger Co.}
\vspace{.1em}
\textbf{Value:} \$5.20 adjusted EPS $\times$ 15 = \$78.00, versus \$67.70 entry. FCFF cross-check: \$87.45.
\begin{itemize}\tight
\item Grocery demand and cash conversion support normalized earnings. The fair multiple discounts heterogeneous retail peers.
\item \textbf{Catalyst: 3 December} results. Watch identical sales, fulfillment margins and guidance.
\item \textbf{Main risk:} leverage and buybacks. Dividends are covered about 2.40 times, but buybacks exceed cash remaining after dividends.
\item \textbf{What would change our view:} a guidance cut or EPS growth driven by debt-funded buybacks rather than operations.
\end{itemize}
\smnote{Fund target 8.10\%. Original annual alpha 8.14\%. Sources: Kroger Q2 FY2026 release and Q3 conference-call announcement.}
\end{frame}


\begin{frame}{HRL -- Hormel Foods}
\vspace{.1em}
\textbf{Value:} \$1.48 adjusted EPS $\times$ 17 = \$25.16, versus \$21.95 entry. FCFF cross-check: \$31.49.
\begin{itemize}\tight
\item We pay for a margin recovery. The 17x multiple is a judgmental premium to food peers, requiring delivery.
\item \textbf{Catalyst: 3 December} results and FY2027 outlook. Watch Retail volumes, Foodservice and cash generation.
\item \textbf{Main risk:} Brakebush funding and thin dividend cover. Coverage falls from 1.17x to about 1.00x in the combined debt/refinancing stress.
\item \textbf{What would change our view:} weak FY2027 guidance or acquisition funding that strains the dividend.
\end{itemize}
\smnote{Fund target 6.90\%. Original annual alpha 7.79\%. Sources: Hormel Q3 release, earnings-call and Brakebush announcements.}
\end{frame}


\begin{frame}{BMY -- Bristol Myers Squibb}
\vspace{.1em}
\textbf{Value:} \$6.45 non-GAAP EPS $\times$ 11 = \$70.95, versus \$62.14 entry. Flat-cash FCFF value: \$65.50.
\begin{itemize}\tight
\item Current medicines generate cash, but patent losses justify a low multiple. Disclosed IPRD costs are retained.
\item \textbf{Catalyst:} expected year-end milvexian/Sotyktu readouts and December scientific updates. Exact readout dates are unconfirmed.
\item \textbf{Main risk:} the cash cushion is thin. A 2\% annual cash decline for five years gives \$58.59; this is a stress, not our flat base case.
\item \textbf{What would change our view:} adverse trial data or growth medicines failing to replace legacy revenues.
\end{itemize}
\smnote{Fund target 6.33\%. Original annual alpha 7.52\%. Source: BMY Q3 FY2026 earnings release and results presentation.}
\end{frame}
\begin{frame}{XOM -- Exxon Mobil}
\vspace{.1em}
\textbf{Value:} our normalized \$10 EPS $\times$ 18 = \$180.00, versus \$165.87 entry. Same-scenario FCFF value: \$213.08.
\begin{itemize}\tight
\item Low-cost production growth supports the normalized case. \$10 EPS is our assumption, not issuer guidance.
\item \textbf{Catalyst:} possible Q4 Uaru start-up and corporate-plan news. No fixed date; this is our weakest timing case.
\item \textbf{Main risk:} the oil cycle and existing DBC exposure. A lower historical-CFO scenario gives \$166.61, near entry.
\item \textbf{What would change our view:} weaker oil/refining margins or project delivery slipping beyond the round.
\end{itemize}
\smnote{Fund target 3.62\%. Original annual alpha 4.51\%. Sources: Exxon Q3 FY2026 release/call and 2023--25 cash-flow statements.}
\end{frame}
\begin{frame}{HPQ -- HP Inc.}
\vspace{.1em}
\textbf{Value:} \$2.38 normalized GAAP EPS $\times$ 13.5 = \$32.13, versus \$28.46 entry. FCFF cross-check: \$34.44.
\begin{itemize}\tight
\item We remove \$0.19 of full-year tariff-refund EPS. A temporary benefit is not recurring earning power.
\item \textbf{Catalyst: 24 November}, confirmed Q4 results and FY2027 outlook. Watch PC margins, memory costs and printing.
\item \textbf{Main risk:} printing decline and cash-funded payouts. Negative equity makes P/B, ROE and debt/equity poor valuation checks.
\item \textbf{What would change our view:} FY2027 guidance below normalized expectations or weaker cash conversion.
\end{itemize}
\smnote{Fund target 3.73\%. Original annual alpha 7.23\%. Sources: HP Q3 FY2026 release and 3 November earnings-call announcement.}
\end{frame}
\begin{frame}{Cash-Flow and Timing Checks}
\vspace{.1em}
\small
\begin{tabular}{lrrrrr}\toprule
Stock & 5-year g & Terminal g & WACC & FCFF value & Review value\\\midrule
ADBE&8.0\%&3.0\%&13.08\%&\$271.14&\$271.14\\
KR&3.0\%&2.0\%&6.31\%&\$87.45&\$78.00\\
HRL&2.0\%&2.0\%&6.20\%&\$31.49&\$25.16\\
BMY&0.0\%&0.0\%&5.83\%&\$65.50&\$65.50\\
XOM&2.0\%&2.0\%&5.83\%&\$213.08&\$180.00\\
HPQ&4.0\%&3.0\%&10.94\%&\$34.44&\$32.13\\
\bottomrule\end{tabular}
\begin{itemize}\tight
\item FCFF and FCFE reconcile under the same financing assumptions. Review value is the lower of cash and earnings values, with weights unchanged.
\item Implied growth is below our base assumptions, but depends on rates, terminal growth and cash normalization.
\item Slower/no rerating is tested with dividend-plus-sale IRRs. A 50\% alpha haircut is not a measured convergence probability.
\end{itemize}
\end{frame}
\begin{frame}{Rejected Candidates}
\vspace{.1em}
\begin{itemize}\tight
\item We researched nine positive-alpha cases and compared 420 subsets containing three to six stocks.
\item FDS, INTU and GILD remain alternatives. They do not justify replacing the agreed six-stock basket.
\item The five-stock full-residual IR maximum is 0.698; our six-stock IR is 0.696.
\item Six stocks reduce recent TE from 8.68\% to 7.74\% and maximum stock weight from 9.09\% to 8.10\%.
\end{itemize}
\vspace{1em}\smnote{A risk/concentration tradeoff within our researched universe. We do not claim a global optimum. The largest position still slightly exceeds our discretionary 8\% preference.}
\end{frame}


% ---------------------------- 3. ALPHAS ----------------------------
\begin{frame}{Alphas and Treynor-Black Weights}
\vspace{.1em}
\small $w_i=\frac{1}{3}\frac{\max(\alpha_i,0)/\sigma_{\varepsilon i}^2}{\sum_j\max(\alpha_j,0)/\sigma_{\varepsilon j}^2}$, with negative-alpha names receiving zero.\par\medskip
\begin{tabular}{lrrrr}\toprule
Stock & Base annual alpha & Residual vol. & SAA beta & Fund weight\\\midrule
ADBE&12.07\%&34.81\%&2.086&4.65\%\\
KR&8.14\%&21.66\%&1.255&8.10\%\\
HRL&7.79\%&22.96\%&0.521&6.90\%\\
BMY&7.52\%&23.55\%&0.989&6.33\%\\
XOM&4.51\%&24.12\%&1.681&3.62\%\\
HPQ&7.23\%&30.08\%&2.122&3.73\%\\
\bottomrule\end{tabular}\par\medskip
\smnote{Original annual alpha = $0.5(1+k)(V/P-1)$. CAPM k uses the market beta; residual risk uses the SAA regression. The 50\% haircut is a common judgment. Quantities use full-precision weights and whole-share rounding.}
\end{frame}


\begin{frame}{Fundamental Law and Credibility}
\vspace{.1em}
\begin{itemize}\tight
\item Classroom diagonal residual IR: \textbf{0.756}. With residual correlations: \textbf{0.696}. Both are model forecasts.
\item Fundamental Law: $IR=IC\sqrt{N}$. With six independent ideas and $IC=0.10$, IR is only 0.245. Our forecast implies much greater skill.
\item The theoretical $S_P^2=S_B^2+IR_A^2$ requires optimal mixing and valid alphas. It does not certify our fixed one-third portfolio.
\item Conservative cash valuations lower fixed-weight IR to \textbf{0.564}. Forward checks also reveal risk underestimation: 80\% of ten monthly-pass origins later breach 5\% daily TE.
\end{itemize}
\smnote{We retain the agreed holdings, with explicit valuation, timing and risk limitations. Historical tests have not established an outperformance advantage.}
\end{frame}


% ---------------------------- 4. PERFORMANCE ----------------------------






% ---------------------------- CLOSING ----------------------------
{
\setbeamertemplate{footline}{}
\setbeamercolor{background canvas}{bg=paper}
\begin{frame}[plain]
\centering
\vfill
\includegraphics[width=5cm]{Meridian-Capital-Official-Logo.png}\par
\vspace{0.4cm}
{\color{gold}\rule{3cm}{0.8pt}}\par
\vspace{0.4cm}
{\rmfamily\fontsize{22}{26}\selectfont\color{navy}Thank you\par}

\vfill
\end{frame}
}

\end{document}
